%
% ellipse.tex -- Ellipse mit hervorgehobenem Viertelbogen
%
% Parametrisierung (a cos t, b sin t) mit a = 2.2, b = 1.4.
% Der erste Quadrant (t in [0, pi/2]) ist hervorgehoben, er entspricht
% dem Integral von 0 bis pi/2 in U = 4a E(k).
% Gleicher Achsenmassstab in x und y.
%
\documentclass[tikz]{standalone}
\usepackage{amsmath}
\usepackage{times}
\usepackage{txfonts}
\usetikzlibrary{arrows,math}
%
% farben.tex -- Farben aus buch/common/macros.tex
%
\definecolor{darkgreen}{rgb}{0,0.6,0}
\definecolor{darkred}{rgb}{0.8,0,0}
\definecolor{orange}{rgb}{1,0.6,0}
\definecolor{gelb}{rgb}{1,0.8,0}

\begin{document}
\def\skala{1.3}
\begin{tikzpicture}[>=latex, thick, scale=\skala]

\def\a{2.2}
\def\b{1.4}

% Achsenlabels verschieben
\def\xdx{-0.05}\def\xdy{-0.05}
\def\ydx{-0.05}\def\ydy{0.0}

% Achsen, gleiche Skalierung in x und y
\draw[->] (-2.9,0) -- (2.9,0);
\draw[->] (0,-2.1) -- (0,2.1);
\node at ({2.9+\xdx},{\xdy}) [anchor=north east] {$x$};
\node at ({\ydx},{1.8+\ydy}) [anchor=south east] {$y$};

% volle Ellipse
\draw (0,0) ellipse ({\a} and {\b});

% erster Quadrant hervorgehoben
\draw[darkred, line width=1.8pt]
	(\a,0) arc[start angle=0, end angle=90, x radius=\a, y radius=\b];

% Halbachsen mit Labels
\draw[thin] (0,0) -- (\a,0) node[midway, below] {$a$};
\draw[thin] (0,0) -- (0,\b) node[midway, left] {$b$};
\fill (\a,0) circle (0.06);
\fill (0,\b) circle (0.06);

\end{tikzpicture}
\end{document}
